\documentclass[11pt,a4paper]{article}
\usepackage[utf8]{inputenc}
\usepackage[T1]{fontenc}
\usepackage[french]{babel}
\usepackage{geometry}
\usepackage{graphicx}
\usepackage{listings}
\usepackage{xcolor}
\usepackage{hyperref}
\usepackage{amssymb}

\geometry{margin=2.5cm}
\lstset{
  basicstyle=\ttfamily\small,
  keywordstyle=\color{blue},
  commentstyle=\color{gray},
  breaklines=true,
  frame=single
}

\title{T2112 - Devoir 1 \\Rapport}
\author{Nom : KETELS Noah}

\begin{document}
\maketitle
\tableofcontents
\bigskip

\section*{Résumé}

L'objectif de ce travail était d'étendre l'application de gestion d'utilisateurs existante en y ajoutant un système de groupes. Cela permet désormais d'organiser les utilisateurs par équipe ou par fonction (ex: "Serveurs", "Cuisine") et de filtrer l'affichage en conséquence côté frontend.

\section{Objectifs (ce que je souhaite atteindre)}

Mon but était d'ajouter une fonctionnalité de gestion de groupes complète et fonctionnelle :
\begin{itemize}
    \item Créer un modèle \texttt{Group} en base de données.
    \item Établir une relation \texttt{User} $\rightarrow$ \texttt{Group} (un utilisateur appartient à un groupe).
    \item Ajouter des endpoints API pour lister et créer des groupes.
    \item Modifier l'endpoint \texttt{GET /api/users} pour qu'il renvoie aussi le groupe de chaque utilisateur et permette le filtrage par \texttt{groupId}.
    \item Mettre à jour l'interface HTML/JavaScript pour permettre la création de groupes et l'affectation d'un groupe lors de l'ajout/modification d'un utilisateur.
\end{itemize}

\section{Stratégie et raisonnement}

J'ai suivi une approche incrémentale pour ne pas tout casser d'un coup :
\begin{enumerate}
    \item \textbf{Côté base de données :} Définir d'abord le modèle \texttt{Group} avec Sequelize, puis ajouter la clé étrangère \texttt{groupId} dans le modèle \texttt{User}. J'ai choisi une relation simple : un groupe peut avoir plusieurs utilisateurs, mais un utilisateur n'a qu'un seul groupe (ou aucun, car la clé est nullable).
    \item \textbf{Côté serveur (API) :} Créer les routes pour les groupes (\texttt{GET /api/groups}, \texttt{POST /api/groups}) et adapter la route des utilisateurs pour inclure le modèle associé (utilisation de l'option \texttt{include} de Sequelize). J'ai aussi ajouté une validation basique des entrées pour éviter les erreurs 500.
    \item \textbf{Côté client :} Charger dynamiquement la liste des groupes dans un \texttt{<select>} au moment du chargement de la page et lors de l'affichage du formulaire. Afficher le nom du groupe sous forme de badge à côté de chaque utilisateur.
\end{enumerate}

\section{Réalisation (ce que j'ai réellement atteint et comment)}
\subsection{Fonctionnalités ajoutées}
\begin{itemize}
    \item \textbf{Modèle Group :} Fichier \texttt{src/models/Group.ts}
    \item \textbf{Mise à jour du modèle User :} Ajout de la propriété \texttt{groupId} et définition de la relation \texttt{belongsTo} dans \texttt{src/models/User.ts}
    \item \textbf{Associations :} Définition des relations croisées dans \texttt{src/models/associations.ts}
    \item \textbf{Nouvelles routes API :} \texttt{src/routes/groups.ts} avec les endpoints \texttt{GET /api/groups} et \texttt{POST /api/groups}
    \item \textbf{Modification de la route users :} Filtrage par \texttt{groupId} via query parameter et inclusion du groupe dans la réponse (\texttt{src/routes/users.ts})
    \item \textbf{Interface utilisateur :} Ajout d'un champ de sélection de groupe dans le formulaire principal et affichage du groupe dans la liste des utilisateurs (\texttt{public/index.html} et \texttt{public/script.js})
\end{itemize}

\subsection{Exemple de code}
J'ai modifié la route \texttt{GET /api/users} pour qu'elle accepte un paramètre \texttt{groupId} dans la query string. Si ce paramètre est présent, on filtre les utilisateurs sur cette valeur. On utilise aussi \texttt{include} pour récupérer les infos du groupe associé.

\begin{lstlisting}[caption={Extrait de src/routes/users.ts}]
router.get('/api/users', async (req, res) => {
  try {
    const { groupId } = req.query;
    const whereCondition = groupId ? { groupId } : {};
    
    const users = await User.findAll({
      where: whereCondition,
      include: [{ model: Group, as: 'group' }]
    });
    
    res.json(users);
  } catch (error) {
    res.status(500).json({ error: 'Erreur serveur' });
  }
});
\end{lstlisting}

Pour la création d'un groupe, j'ai ajouté un endpoint simple avec validation :

\begin{lstlisting}[caption={Extrait de src/routes/groups.ts}]
router.post('/api/groups', async (req, res) => {
  try {
    const { name } = req.body;
    
    if (!name || name.trim() === '') {
      return res.status(400).json({ 
        error: 'Le nom du groupe est requis' 
      });
    }
    
    const newGroup = await Group.create({ 
      name: name.trim() 
    });
    res.status(201).json(newGroup);
  } catch (error) {
    res.status(500).json({ error: 'Erreur serveur' });
  }
});
\end{lstlisting}

\subsection{Captures d'écran}
Voici le résultat obtenu après l'implémentation. On voit bien le nouveau champ "Groupe" dans le formulaire et les badges à côté de chaque utilisateur.

\begin{figure}[h]
  \centering
  \includegraphics[width=0.9\textwidth]{screenshots/liste-utilisateurs-avec-groupes.png}
  \caption{Capture d'écran de l'interface avec la liste des utilisateurs et leurs groupes respectifs}
  \label{fig:interface-groupes}
\end{figure}

\begin{figure}[h]
  \centering
  \includegraphics[width=0.9\textwidth]{screenshots/creation-groupe.png}
  \caption{Formulaire de création d'un nouveau groupe}
  \label{fig:creation-groupe}
\end{figure}

\newpage
\section{Annexes}
\subsection{Code source - src/models/Group.ts}
\label{ann:group-model}
\begin{lstlisting}
import { DataTypes, Model } from 'sequelize';
import sequelize from '../config/database.js';

class Group extends Model {
  declare id: number;
  declare name: string;
}

Group.init({
  id: {
    type: DataTypes.INTEGER,
    autoIncrement: true,
    primaryKey: true
  },
  name: {
    type: DataTypes.STRING,
    allowNull: false
  }
}, {
  sequelize,
  tableName: 'groups',
  timestamps: true
});

export default Group;
\end{lstlisting}

\subsection{Code source - Extrait de src/models/associations.ts}
\label{ann:associations}
\begin{lstlisting}
import Group from './Group.js';
import User from './User.js';

Group.hasMany(User, { 
  foreignKey: 'groupId', 
  as: 'users' 
});

User.belongsTo(Group, { 
  foreignKey: 'groupId', 
  as: 'group' 
});

export { User, Group };
\end{lstlisting}

\end{document}